À $300$ K, le calcul donne $u=-3.813000$ eV/atome et $q=1.0000$. La comparaison à la borne exacte vérifie l'état ordonné de basse température.

Au maximum thermique, à $1050$ K, $q=0.575$ et $c_V/k_B=5.87$ par atome. Les historiques et les distributions permettent d'évaluer l'exploration des configurations dans cette zone.

À $4000$ K, on obtient $u=-3.757820$ eV/atome, $q=0.065$, $\alpha_1=-0.067$ et $\alpha_2=0.026$. Ces valeurs sont celles d'une température finie ; elles ne doivent pas être remplacées par les limites de mélange indépendant.

Les critères locaux fixés avant l'analyse sont satisfaits dans la fenêtre du maximum. Cela appuie la cohérence statistique du résultat à cette taille, tout en laissant ouvertes les erreurs systématiques liées à des états métastables non visités.
